\FloatBarrier
\section{Converting Spin States into Readout Probabilities}
\label{app:spin-readout}
\label{sec:readout-calibration}

Spin readout requires the possible spin states to produce distinguishable
charge records. This appendix develops the conversion, timing, and
assignment of those records for the encodings in
Section~\ref{sec:initialization-readout}, then relates repeated outcomes
to logical-state probabilities. For an array, the sensor responses
must also distinguish which qubits produced the charge signals.
The final section develops the separation of several outcomes and
the recovery of their joint probabilities, supporting the readout
requirements in Section~\ref{sec:platform-array-selection}.

The sensor is a separate conducting device whose current changes
when nearby target electrons move. For a sensor dot, the Coulomb
term $\sum_i U_{i\mathrm s}N_i n_{\mathrm s}$ shifts its addition
energy when the target occupations $N_i$ change. At fixed sensor
source and drain voltages, the local current response is
\begin{equation}
 \delta I_{\mathrm s}\simeq
 \frac{\partial I_{\mathrm s}}{\partial\mu_{\mathrm s}}
             \sum_i U_{i\mathrm s}\delta N_i .
 \label{eq:calibration-sensor-charge-response}
\end{equation}
Here $n_{\mathrm s}$ is the sensor occupation, $\mu_{\mathrm s}$
its addition energy and $U_{i\mathrm s}$ the interaction energy
with target dot $i$. The sensor is placed on a sloping part of its
current curve, as in Section~\ref{sec:measurement-charge-sensor}.

To read a spin, we must first make the two possible spin states
produce different target charge records. An electron can leave
for one single-spin state, or a pair can enter a doubly occupied
dot only when it is a singlet. We choose the voltage and duration
of that conversion, then infer the spin outcome from the
resulting sensor trace.

\subsection{A single-spin qubit}
\label{sec:calibration-single-spin-readout}

The spin state is initially stored in an occupied dot isolated from
its reservoir. To read it, raise the dot potential and reopen the
reservoir barrier so that an up electron can leave but a down
electron cannot. With the notation of
Section~\ref{sec:calibration-single-spin-initialization}, the
readout setting $\mathbf V^{\mathrm R}$ must satisfy
\begin{equation}
 \mu_\downarrow=\epsilon_d-\frac{E_Z}{2}
   <\mu_{\mathrm R}<
 \mu_\uparrow=\epsilon_d+\frac{E_Z}{2},
 \qquad
 \mu_{\mathrm R}-\frac{E_Z}{2}<\epsilon_d<
       \mu_{\mathrm R}+\frac{E_Z}{2}.
 \label{eq:calibration-readout-window}
\end{equation}
Keep the second-electron addition threshold above the reservoir.
The plunger sets $\epsilon_d$ and the barrier sets the tunnel
rates; use their joint voltage response to reach this condition.

At zero temperature, the reservoir has empty states at
$\mu_\uparrow$ but occupied states at $\mu_\downarrow$.
An up electron can therefore escape. Once the dot is empty, a
down electron can enter, giving the charge sequence
\[
 \uparrow:\quad N=1\longrightarrow0\longrightarrow1,
 \qquad
 \downarrow:\quad N=1.
\]
The information is the temporary empty interval, not the final
electron number. A nearby sensor has different mean signals
$s_0$ and $s_1$ for an empty and occupied target. A simple record
model is
\begin{equation}
 s(t)=s_0+(s_1-s_0)N(t)+\xi(t),
 \label{eq:calibration-single-spin-record}
\end{equation}
where $\xi(t)$ is detector noise. Thus the up-spin record can
contain an excursion to $s_0$, while the down-spin record stays
near $s_1$. Resolved empty and occupied intervals in charge
reference traces at these readout voltages determine $s_0$ and
$s_1$; using the same voltages also fixes the direct gate response
of the sensor.

To locate the readout interval, measure the two first-electron
thresholds relative to $\mu_{\mathrm R}$, then compare records
after down-spin preparation and after a spin-flip pulse that
prepares up. Vary the readout plunger voltage while keeping the
reservoir contact fixed. Choose a setting between the thresholds,
with margins against thermal activation. This window distinguishes
an already present spin through its ability to leave~\cite{Elzerman2004}.

The basic readout is now specified: move the isolated electron to
$\mathbf V^{\mathrm R}$, allow spin-selective tunneling for a time
$t_r$, and record whether a temporary empty interval occurs.
An excursion reports up; no excursion reports down in the ideal
conversion. The finite-time and relaxation errors in that rule
determine the choice of $t_r$ in
Section~\ref{sec:calibration-readout-duration-choice}.

\subsection{A singlet--triplet qubit}
\label{sec:calibration-st-readout}

The two logical states $S_{11}$ and $T_{0,11}$ initially have the
same charge configuration. To distinguish them, lower the energy
of double occupation on dot 2 while keeping the reservoirs closed.
Spin-conserving tunneling connects $S_{11}$ to $S_{02}$. A triplet
cannot place both electrons in the same lowest orbital--valley
state, because its symmetric spin state requires an antisymmetric
spatial state. Its lowest doubly occupied state is higher in
energy. Choose the detuning so that
\begin{equation}
 E(S_{02})<E(S_{11}),\qquad
 E(T_{02,m})>E(T_{0,11})
 \quad\hbox{for relevant triplet final states}.
 \label{eq:calibration-pair-readout-window}
\end{equation}
Both inequalities refer to the same readout voltages. Together
with small spin mixing during conversion, they allow
$S_{11}\to S_{02}$ while keeping $T_{0,11}$ in $(1,1)$.
This is the spin-dependent obstruction to charge transfer called
Pauli spin blockade.

The voltage pulse reverses the singlet-separation ramp of
Section~\ref{sec:calibration-st-initialization}. Follow the
singlet charge anticrossing into $(0,2)$ without waiting long
enough for substantial singlet--triplet mixing or relaxation.
At the endpoint, a sensor with local response
$s=s_{\mathrm{bg}}+\kappa_1N_1+\kappa_2N_2+\xi$
has mean levels
\[
 s_{11}=s_{\mathrm{bg}}+\kappa_1+\kappa_2,\qquad
 s_{02}=s_{\mathrm{bg}}+2\kappa_2,\qquad
 s_{02}-s_{11}=\kappa_2-\kappa_1.
\]
The coefficients $\kappa_i$ are signal changes per electron on
dot $i$. The sensor must respond differently to the two dots.
If the converted state persists, averaging the plateau signal
distinguishes singlet from triplet.

To set the readout voltage, locate the $(1,1)$--$(0,2)$ boundary
and scan across it using singlet and triplet reference preparations.
The singlet reference comes from double-occupation loading.
For a $T_0$ reference, apply the gradient-induced $x$ rotation of
Eq.~\eqref{eq:st-axis-limits} to a separated singlet: when exchange
is negligible, a wait $\pi\hbar/|\Delta E_Z|$ interchanges $S$ and
$T_0$. Determine $|\Delta E_Z|$ from the period of
$P_{T_0}(t)=\sin^2(\Delta E_Zt/2\hbar)$ while varying this wait;
the charge-conversion contrast sets the observed oscillation amplitude.
Compare the two charge-transfer probabilities as functions of the
readout voltage. Choose the interval in which the singlet transfers
but the triplet remains blocked. For the coherent conversion used
here, select the ramp speed from
the charge-following and spin-preservation conditions in
Section~\ref{sec:calibration-st-initialization}. After the ramp,
record the charge plateau before appreciable triplet unblocking.

The resulting map is $S_{11}\mapsto(0,2)$ and
$T_{0,11}\mapsto(1,1)$. The ramp performs the conversion; the
subsequent plateau supplies the sensor record. A hold that instead
relies on irreversible singlet transfer is a different implementation,
treated separately below when choosing its readout duration.

\subsection{An exchange-only qubit}
\label{sec:calibration-eo-readout}

Read the spin of the pair on dots 1--2 while dot 3 remains occupied.
For logical 0, the pair is a singlet and can undergo
$(1,1,1)\rightarrow(0,2,1)$; for logical 1, it is a triplet and
remains blocked in $(1,1,1)$. Suppress $J_{23}$ during conversion
so that exchange with dot 3 does not change the pair being measured.
Use the pair-readout energy window of
Eq.~\eqref{eq:calibration-pair-readout-window}, adding the fixed
occupation of dot 3 to each charge state and measuring the energies
with that electron present. The sensor contrast is now
$s_{021}-s_{111}=\kappa_2-\kappa_1$; the unchanged third-dot
contribution cancels. The same voltage and duration tests then
apply. Within the logical space, pair singlet and pair triplet
are the two outcomes even though logical 1 need not be a product
of a definite spin on dot 3 and a definite triplet component.

This readout distinguishes pair singlet from pair triplet, but
does not determine the third spin. In particular, every state
in Eq.~\eqref{eq:calibration-eo-preparation-state} produces the
same pair-singlet result, regardless of $p_3$. Testing preparation
in the fixed-$m=+1/2$ encoding therefore also requires the separate
polarization measurement described in
Section~\ref{sec:calibration-eo-initialization}.

\subsection{Choosing a duration that preserves the spin contrast}
\label{sec:calibration-readout-duration-choice}

The conversion conditions allow the desired charge motion, but do
not guarantee that it happens before recording ends. Waiting longer
can recover slow desired events while also admitting unwanted ones.
We now hold the readout voltages fixed and choose the recording
time from the probabilities of these competing events.

\subsubsection{A single spin: escape before relaxation}
For an up spin, the useful event is escape to the reservoir.
Relaxation to down before escape destroys that indication.
For a down spin, thermally activated escape gives a false indication.
The time should allow the first process and suppress the other two.

For a reservoir at temperature $T$, a weak spin-conserving tunnel
contact with bare rate $\Gamma_s^{(0)}$ for spin
$s\in\{\uparrow,\downarrow\}$ gives
\begin{equation}
 \Gamma_{\mathrm{out},s}
 =\Gamma_s^{(0)}[1-f(\mu_s)],\qquad
 f(E)=\frac{1}{1+e^{(E-\mu_{\mathrm R})/(k_{\mathrm B}T)}}.
 \label{eq:calibration-spin-exit-rates}
\end{equation}
Raising both dot addition energies increases the desired up-spin
exit rate, but eventually also allows a down spin to escape.
The voltage scan therefore selects the rates, not just the ordering
of three energies.

The readout duration must allow the desired exit before spin
relaxation. Write $\Gamma_u=\Gamma_{\mathrm{out},\uparrow}$ and
$\Gamma_d=\Gamma_{\mathrm{out},\downarrow}$.
At $E_Z\gg k_BT$, upward spin excitation is negligible and the
downward relaxation rate is $\gamma_1\simeq T_1^{-1}$ at the
readout setting. An initial up spin survives without either
exit or relaxation until time $u$ with probability
$e^{-(\Gamma_u+\gamma_1)u}$. Its desired exit probability is
therefore
\begin{equation}
 \begin{aligned}
 P_{\mathrm{exit}}(t_r)
 &=\int_0^{t_r}\Gamma_u e^{-(\Gamma_u+\gamma_1)u}\,du\\
 &=\frac{\Gamma_u}{\Gamma_u+\gamma_1}
          \left[1-e^{-(\Gamma_u+\gamma_1)t_r}\right].
 \end{aligned}
 \label{eq:calibration-readout-duration}
\end{equation}
This is the competing-rate model of
Eq.~\eqref{eq:readout-competing-rates} for the present voltage
setting. To choose a duration,
consider the regime $\Gamma_d t_r\ll1$, in which an up spin that
relaxes before escaping is normally missed. A down-spin reference
produces a false exit with probability $1-e^{-\Gamma_d t_r}$,
whereas the desired up-spin exit probability is the expression
above. For equally likely input spins, the conversion error in
this approximation is
\begin{equation}
 \epsilon_{\mathrm{conv}}(t_r)
 =\frac12\left[
  1-\frac{\Gamma_u}{k}(1-e^{-kt_r})
  +1-e^{-\Gamma_d t_r}\right],
 \qquad k=\Gamma_u+\gamma_1.
 \label{eq:calibration-conversion-error}
\end{equation}
The first contribution decreases with time; the second increases.
For $\Gamma_u>\Gamma_d>0$, their balance gives
\begin{equation}
 \frac{d\epsilon_{\mathrm{conv}}}{dt_r}=0
 \quad\Longrightarrow\quad
 t_r^\star=\frac{\ln(\Gamma_u/\Gamma_d)}{k-\Gamma_d}.
 \label{eq:calibration-optimal-readout-time}
\end{equation}
Check $\Gamma_d t_r^\star\ll1$ before using this estimate. For
example, the illustrative rates
$(\Gamma_u,\Gamma_d,\gamma_1)=(10^4,10,100)\,\mathrm{s}^{-1}$
give $t_r^\star\simeq0.685\,\mathrm{ms}$ and about $0.89\%$
conversion error. Increasing the wait alone does not remove the
competition with relaxation.

Repeated up/down reference traces determine the exit probabilities
versus time. A separate relaxation measurement fixes $\gamma_1$;
the up-spin curve then determines $\Gamma_u$, and the down-spin
curve determines $\Gamma_d$. Repeat at nearby readout voltages and
barrier settings, and compare the resulting minimum errors.
Finally test the same settings with the actual sensor decision rule.
These formulas assume that an exit event is resolved; an electron
that leaves and refills too quickly to be detected adds a separate
missed-event contribution.

\subsubsection{A pair readout implemented by irreversible transfer}
In the coherent protocol above, a ramp carries $S_{11}$ into $S_{02}$.
An alternative is to move to a fixed detuning where $S_{02}$ is lower,
then let tunneling and energy dissipation complete the transfer.
An initially blocked triplet may also change spin and transfer later.
The two inputs consequently produce different charge-transfer
probabilities at intermediate times.

Here the clock starts with the pair still in $(1,1)$ at the conversion
setting. This is not a rate description of the preceding coherent ramp.
When the fixed-setting records support an incoherent, one-way model,
write
\[
 S_{11}\xrightarrow{\Gamma_S} S_{02},\qquad
 T_{0,11}\xrightarrow{\gamma_T}S_{11}
                        \xrightarrow{\Gamma_S}S_{02}.
\]
Here $\Gamma_S$ is the singlet transfer rate including energy
dissipation, and $\gamma_T$ describes spin-changing processes that
unblock a triplet. Assume that the converted charge remains
detectable and that reverse processes are negligible at this
endpoint. Starting from $S$ gives
\[
 p_{02|S}(t)=1-e^{-\Gamma_S t}.
\]
Starting from $T_0$, the state must first unblock. Integrating over
its unblocking time gives, for $\Gamma_S\ne\gamma_T$,
\begin{align}
 p_{02|T_0}(t)
 &=\int_0^t\gamma_Te^{-\gamma_T u}
                         [1-e^{-\Gamma_S(t-u)}]\,du \notag\\
 &=1-\frac{\Gamma_S e^{-\gamma_Tt}
                 -\gamma_Te^{-\Gamma_St}}{\Gamma_S-\gamma_T}.
 \label{eq:calibration-blockade-conversion}
\end{align}
A singlet that has not yet transferred and a triplet that has already
unblocked both give a wrong charge label. For equally likely inputs,
their conversion error at a chosen time is
\[
 \epsilon_{\mathrm{pair}}(t)
 =\frac12[1-p_{02|S}(t)+p_{02|T_0}(t)].
\]
Thus choosing $t$ to maximize $p_{02|S}-p_{02|T_0}$ minimizes this
error. That difference is large in the interval
$\Gamma_S^{-1}\ll t\ll\gamma_T^{-1}$,
when such an interval exists. Reference time scans determine
whether this separation of rates is actually available.
The fitted rates characterize this particular conversion setting.
Moving the voltage or changing the ramp requires repeating the
conversion test.

For either pair-conversion implementation, the subsequent sensor
average must resolve the two levels before the charge contrast is
lost. Its duration is tested on the actual traces. A measured
irreversible transfer rate cannot be substituted for the coherent
ramp condition.

\subsection{Assigning 0 or 1 to the measured record}
\label{sec:calibration-record-classification}

The two charge configurations produce different mean sensor levels,
but noise makes the observed records overlap. We must specify
which feature of a record counts as logical 0 or 1.
The conversion and its duration are already fixed. We now determine
how to extract an outcome from its noisy sensor trace; this step
does not change the physical spin-to-charge mapping.

\subsubsection{Choosing the feature to classify}
For pair readout, average over a fixed plateau interval of length
$t_{\mathrm a}$ to obtain $\bar s$. The normalized statistic
\[
 x=\frac{\bar s-s_{02}}{s_{11}-s_{02}}
\]
is near zero for a converted singlet and near one for a blocked
triplet. For the three-dot encoding, use $s_{021}$ and $s_{111}$.

Single-spin readout requires a different statistic because the
empty interval is temporary. After averaging samples over the
detector response time, one simple choice is
\[
 x=\max_{0<t<t_r}\frac{s(t)-s_1}{s_0-s_1}.
\]
An empty-dot excursion moves this statistic toward one; an
occupied-dot record stays near zero. The reference traces must
use the same sampling, averaging and readout duration, because
all three affect the chance of a noise excursion or a missed
short event. For either statistic, a threshold $x_c$ reports
logical 1 when $x>x_c$.

\subsubsection{Choosing a decision threshold}
Reference preparations of 0 and 1 give two distributions of $x$.
Count their misclassified records:
\[
 \epsilon_0=P(x>x_c\mid0),\qquad
 \epsilon_1=P(x\le x_c\mid1).
\]
For equally weighted reference states, choose $x_c$ to minimize
$(\epsilon_0+\epsilon_1)/2$. Keep the voltage pulse, record duration
and decision rule the same in the subsequent experiment.

The threshold can be selected directly from the two measured
distributions. To see analytically how their separation controls
the error, suppose a plateau-averaged signal has Gaussian
conditional distributions with means $m_0<m_1$ and the same
standard deviation $\sigma$. At a proposed threshold,
\[
 \epsilon_0=\int_{x_c}^{\infty}
  \frac{e^{-(x-m_0)^2/(2\sigma^2)}}{\sqrt{2\pi}\sigma}\,dx,
 \qquad
 \epsilon_1=\int_{-\infty}^{x_c}
  \frac{e^{-(x-m_1)^2/(2\sigma^2)}}{\sqrt{2\pi}\sigma}\,dx.
\]
Differentiating their sum equates the two probability densities,
giving $x_c=(m_0+m_1)/2$. The error for either state is then
\[
 \epsilon_0=\epsilon_1
 =\frac12\operatorname{erfc}
       \left(\frac{m_1-m_0}{2\sqrt2\,\sigma}\right).
\]
Here $\operatorname{erfc}(z)=(2/\sqrt\pi)\int_z^\infty e^{-u^2}du$.
This example connects histogram separation to assignment error.
For non-Gaussian records, the same threshold scan can be performed
directly by counting the two tails; the Gaussian model is not needed.

\subsubsection{Recovering an outcome probability from repeated records}
The decision rule produces a count of 1 outcomes. If its conditional
errors are independently known, a state with
true logical-1 probability $p$ produces the reported probability
\begin{equation}
 p_{\mathrm{obs}}
 =\epsilon_0(1-p)+(1-\epsilon_1)p,\qquad
 p=\frac{p_{\mathrm{obs}}-\epsilon_0}
         {1-\epsilon_0-\epsilon_1}.
 \label{eq:calibration-assignment-map}
\end{equation}
The second relation requires nonzero readout contrast.
If $n$ repeated trials give the fraction $\hat p_{\mathrm{obs}}$
of outcome 1, and the assignment probabilities are treated as known,
the variance of the corrected estimator is
\begin{equation}
 \operatorname{Var}(\hat p)
 =\frac{p_{\mathrm{obs}}(1-p_{\mathrm{obs}})}
        {n(1-\epsilon_0-\epsilon_1)^2}.
 \label{eq:calibration-assignment-variance}
\end{equation}
The inverse map removes a known bias but amplifies statistical
uncertainty when the two records are difficult to distinguish.
It cannot recover information from zero-contrast readout.
An uncertain reference preparation gives a mixture of the two
distributions, so its histogram alone does not separately determine
preparation error and readout error. Initial calibration can use
well-understood loading and relaxation limits; precision claims
require the reference-state errors to be established as well.
This distinction matters when deciding what a measured gate
probability actually verifies.

\subsection{From repeated readout to a spin probability}
\label{sec:measurement-pulse-cycle}

A single readout assigns one outcome to one preparation. To
measure an outcome probability after a control pulse, repeat the
whole sequence: prepare the same state, isolate the electrons,
apply the same control, and then perform the same charge conversion
and sensor measurement. Reading one spin repeatedly would instead
probe a state already changed by the first readout.

Let $\rho_0$ be the prepared spin density matrix and $U(\tau)$
the coherent evolution during a control interval of duration $\tau$.
For ideal readout with outcome projector $M_a$,
\begin{equation}
 P_a(\tau)=\operatorname{Tr}\!\left[
       M_a U(\tau)\rho_0U^\dagger(\tau)\right].
 \label{eq:measurement-pulse-born-rule}
\end{equation}
The spin-to-charge conversion fixes the measured projector.
For single-spin readout,
$M_\uparrow=\ket{\uparrow}\bra{\uparrow}$ and
$M_\downarrow=\ket{\downarrow}\bra{\downarrow}$.
For singlet--triplet readout, the converted outcome has
$M_S=\ket S\bra S$ within the logical space.
Pair readout of an exchange-only qubit similarly distinguishes
its logical pair-singlet and pair-triplet states, but does not
separately identify every leakage state.

A common timing signal starts the voltage sequence and sensor
recording. For each chosen $\tau$, classify the records from
$K$ independent trials. The fraction $k_a/K$ of assigned outcomes
estimates the observed probability; the assignment map in
Eq.~\eqref{eq:calibration-assignment-map} connects it to $P_a$.
Keep the preparation and readout settings fixed while changing
the intervening control. Its Hamiltonian and propagator determine
how the measured probability varies.

\subsection{Separating several readout outcomes}
\label{sec:processor-array-readout}

A sensor records a charge-dependent signal, not a qubit label.
For sensor \(\ell\), a linear model for its mean signal is
\begin{equation}
 \overline s_\ell=s_{\ell,0}+\sum_i c_{\ell i}\langle N_i\rangle,
 \label{eq:processor-sensor-array}
\end{equation}
where \(N_i\) is the occupation variable after the chosen
spin-to-charge conversion and \(c_{\ell i}\) is the sensor
response to it. If one sensor responds equally to two dots,
it measures their total charge but cannot tell \(10\) from
\(01\). A second sensor, unequal response, or a conversion
at a different time can supply that distinction.

For example, subtract the backgrounds from two sensor
signals and suppose their response is symmetric:
\[
 \begin{pmatrix}\bar s_1\\\bar s_2\end{pmatrix}
 =\begin{pmatrix}c&d\\d&c\end{pmatrix}
   \begin{pmatrix}\langle N_1\rangle\\\langle N_2\rangle\end{pmatrix}.
\]
The sum and difference of the occupations are
\begin{equation}
 \langle N_1\rangle+\langle N_2\rangle
       =\frac{\bar s_1+\bar s_2}{c+d},\qquad
 \langle N_1\rangle-\langle N_2\rangle
       =\frac{\bar s_1-\bar s_2}{c-d}.
 \label{eq:processor-sensor-inversion}
\end{equation}
This inversion needs \(c\ne\pm d\). When \(c\) is close to
\(d\), the difference estimate strongly amplifies signal
noise. An array thus needs sufficiently distinct response
patterns, not simply a count of sensors equal to the count
of qubits.

These equations determine mean occupations, not the joint
probabilities of single-shot outcomes. For example, a distribution
with \(P_{00}=P_{11}=1/2\) and one with
\(P_{01}=P_{10}=1/2\) both have
\(\langle N_1\rangle=\langle N_2\rangle=1/2\), but the first
always has even parity and the second always has odd parity.
Their identical marginal occupations therefore do not determine
their different correlations. To estimate
\(P_{00},P_{01},P_{10},P_{11}\), first
classify the pair of outcomes in each run, then average those
classifications. One sensor can in principle classify all
four if its levels \(s_0,s_0+c_1,s_0+c_2,s_0+c_1+c_2\)
are distinguishable. If \(c_1=c_2\), its \(01\) and \(10\)
levels coincide. Sequential conversion can separate them
in time, provided the later qubit retains its spin state
until its readout.

Readout capacity also sets a time cost. If \(N_m\) qubits
must be measured, at most \(R_m\) independent results can be
obtained at once, and one readout group takes \(t_r\), then
\[
 T_{\mathrm{meas}}\ge
 \left\lceil\frac{N_m}{R_m}\right\rceil t_r.
\]
Conversion, reset, and shared-control delays may increase
this bound. Qubits awaiting measurement must retain their
state for those extra rounds.
